\documentclass[11pt,a4paper]{article}
\usepackage[T1]{fontenc}
\usepackage[utf8]{inputenc}
\usepackage[french]{babel}
\usepackage{geometry}
\usepackage{booktabs}
\usepackage{hyperref}

\geometry{margin=2.2cm}
\setlength{\emergencystretch}{2em}
\hypersetup{colorlinks=true,linkcolor=blue,urlcolor=blue}

\title{Rapport technique, Lab 4\\
\large Cellule NR, grille de 16 terminaux\\
Dépôt \texttt{ns3\_IMT2030\_connectivity}}
\author{Traorehub, travail de recherche en cours}
\date{Novembre 2026}

\begin{document}
\maketitle

\section{Ce qui a été exécuté}

Le module 5G-LENA v3.3.1, étiquette prévue pour ns-3.42, commit \texttt{f13da239}, a été compilé dans l'arbre ns-3.43. Le scénario reprend la grille, le budget et les trois politiques des Labs 3 et 5. La radio est une cellule, à 3,5~GHz, sur 20~MHz, en numérologie 0, en intérieur et en visibilité, sans ombrage. Le planificateur est OFDMA round-robin. Le motif TDD est \texttt{DL|S|UL|UL|UL|DL|S|UL|UL|UL|}. La puissance du gNB est de 30~dBm, celle du terminal de 23~dBm. La capacité de 3~Mbit/s reste le paramètre de la loi $K$. Elle ne règle pas la bande.

La campagne effectivement écrite dans le CSV compte trois runs, tous à 16 terminaux, 40~ms, 10~s et 40 paquets de 512 octets. Le plan complet de 36 runs n'a pas été lancé. Les essais à 36 et à 100 terminaux s'interrompent avant le trafic.

\section{Grille de 16 terminaux}

\begin{center}
\begin{tabular}{@{}lrrrrr@{}}
\toprule
Politique & Octets reçus & Taux & Jain & Événements & Plafond \\
\midrule
baseline & 327\,680 & 1 & 1 & 0 & 16 \\
adaptive & 327\,680 & 1 & 1 & 0 & 16 \\
feedback & 327\,680 & 1 & 1 & 4 & 16 \\
\bottomrule
\end{tabular}
\end{center}

Les 640 paquets prévus arrivent, soit 16 terminaux fois 40 paquets. Le plafond de la loi vaut 24 à 40~ms, puis il est borné par les 16 terminaux. L'admission fixe ne retire donc personne. La boucle enregistre 4 événements et revient à 16, sans changer le volume. Le même volume, 327\,680 octets, est celui du proxy Wi-Fi sur cette cellule non saturée. À cette charge, changer la couche radio ne change pas le service.

\section{Grilles qui ne passent pas le RRC}

À 36 terminaux, puis à 100 terminaux, le processus s'arrête à 30~ms. L'assertion du RRC idéal indique qu'aucun contexte n'est trouvé pour le RNTI 0, dans \texttt{nr-rrc-protocol-ideal.cc} à la ligne 252. Aucun octet UDP n'est alors compté. Ces deux tailles ne figurent pas dans le CSV. Le défaut appartient à cette combinaison module v3.3.1 et ns-3.43. Il n'a pas été contourné par un chiffre inventé.

\section{Lecture}

La cellule NR de 20~MHz absorbe l'offre de 16 terminaux. Elle ne fournit pas encore la comparaison de charge faite sur le Wi-Fi à 64 et 100 terminaux, parce que ces effectifs n'atteignent pas l'émission. La densité de la grille de 16 terminaux reste d'environ $7\times 10^{4}$ terminaux par kilomètre carré. Elle demeure très en dessous d'un million par kilomètre carré.

\section{Reproduction}

Le script \texttt{labs/lab04\_5g\_lena/scripts/lancer\_campagne.py} refait les trois runs de 16 terminaux. Il faut le module \texttt{contrib/nr} à l'étiquette v3.3.1, puis la cible \texttt{imt2030-nr-gestion}. Le CSV est \texttt{labs/lab04\_5g\_lena/sorties/resultats.csv}.

\end{document}
